\chapter{Anatomia della funzione obiettivo}
\label{ch:obiettivo}

Questo capitolo esamina uno per uno i termini della funzione obiettivo: che cosa misurano, perché
sono costruiti in quel modo, che cosa significa concretamente il valore numerico del loro
coefficiente e come cambia il comportamento del modello al variare di quel valore.

\section{Il principio unificante: omogeneità dimensionale}
\label{sec:omogeneita}

La funzione obiettivo somma quantità di natura molto diversa: un errore di ricostruzione, un
conteggio di accensioni, uno scarto di durata, un'eccedenza settimanale. Perché la somma abbia
senso, tutti i termini devono essere \textbf{omogenei}, cioè avere la stessa dimensione fisica.

Il primo termine, l'errore di ricostruzione, si misura in watt. Ne segue che ogni altro termine
deve essere anch'esso in watt. Questo determina la struttura di tutti i coefficienti:

\begin{equation}
  \underbrace{\lambda}_{\text{adimensionale}}
  \;\times\;
  \underbrace{p_i}_{[\mathrm{W}]}
  \;\times\;
  \underbrace{(\text{variabile})}_{\text{adimensionale}}
  \;=\;
  [\mathrm{W}] .
\end{equation}

\begin{notebox}{Perché ogni penalità è moltiplicata per $p_i$}
Non si tratta di una convenzione arbitraria ma di una necessità dimensionale. Se le penalità
fossero costanti assolute, il loro peso relativo dipenderebbe dalla potenza dell'apparecchio: la
stessa penalità sarebbe trascurabile per una lavatrice da $2000\unit{W}$ e proibitiva per un
televisore da $100\unit{W}$. Scalando per $p_i$ ogni apparecchio riceve una penalità
\textbf{proporzionata alla propria scala}, e un unico coefficiente $\lambda$ ha lo stesso
significato per tutti.
\end{notebox}

Da questa costruzione discende l'interpretazione uniforme dei coefficienti, valida per tutti i
termini tranne quello di quota settimanale (§\ref{sec:t-quota}):

\begin{center}
\begin{minipage}{0.9\textwidth}
\centering\itshape
$\lambda = 1$ è il punto di pareggio: violare la dichiarazione una volta costa esattamente quanto
lasciare $p_i$ watt non spiegati per la durata di uno slot.
\end{minipage}
\end{center}

Sotto 1 l'evidenza del segnale prevale sulla dichiarazione dell'utente; sopra 1 prevale la
dichiarazione.

\section{Termine 1 --- Errore di ricostruzione}

\begin{equation}
  \sum_{t \in \Vset} \bigl( \ep_t + \en_t \bigr)
\end{equation}

\paragraph{Che cosa misura.} La distanza complessiva fra il consumo ricostruito dal modello e
quello effettivamente misurato, sommata su tutti gli slot validi della giornata. Come mostrato nel
capitolo~\ref{ch:perche-l1}, per costruzione $\ep_t + \en_t = |\hat{y}_t - y_t|$, quindi il termine
vale l'errore assoluto totale, in watt.

\paragraph{Ruolo nel modello.} È l'unico termine che spinge il modello a \emph{usare} gli
apparecchi: tutti gli altri sono penalità che spingono a non usarli. L'equilibrio fra questo
termine e gli altri determina la soluzione.

\paragraph{Coefficiente.} È fissato a 1 e non è configurabile: costituisce il \textbf{riferimento}
rispetto al quale tutti gli altri coefficienti sono espressi. Cambiarlo equivarrebbe a riscalare
tutti gli altri in senso inverso.

\paragraph{Asimmetria fra sovrastima e sottostima.} Le due variabili hanno lo stesso costo
unitario, quindi il modello è indifferente fra sovrastimare e sottostimare di una stessa quantità.
Questa è una scelta implicita che vale la pena rendere esplicita: se si volesse penalizzare di più
la sovrastima --- attribuire consumo che non esiste --- basterebbe assegnare a $\ep_t$ un
coefficiente maggiore di quello di $\en_t$.

\section{Termine 2 --- Fascia oraria}
\label{sec:t-fascia}

\begin{equation}
  \sum_{i} \sum_{t \in \Vset} \lambda_w \, p_i \, f_i(t) \, x_{i,t}
\end{equation}

\paragraph{Che cosa misura.} Il costo di attribuire consumo a un apparecchio in un momento della
giornata incompatibile con quanto dichiarato nel questionario. È il termine più articolato, perché
il fattore $f_i(t)$ non è binario ma cresce con la distanza dalla fascia dichiarata; la sua
costruzione completa è nel capitolo~\ref{ch:fattore}.

\paragraph{Perché è applicato a $x_{i,t}$ e non a $u_{i,t}$.} La penalità grava su \emph{ogni
slot} in cui l'apparecchio risulta acceso, non solo sull'istante di accensione. Un ciclo di due ore
interamente fuori orario paga quindi otto volte, un ciclo che sconfina di un solo slot paga una
volta. La penalità è cioè proporzionale a \emph{quanto} consumo viene attribuito fuori fascia, non
al semplice fatto che sia avvenuto.

\paragraph{Effetto del coefficiente.}
\begin{center}
\begin{tabularx}{\textwidth}{@{}l X@{}}
\toprule
$\lambda_w = 0$ & la fascia oraria dichiarata viene ignorata del tutto \\
$\lambda_w \approx 0.2$--$0.4$ & la fascia orienta la scelta senza impedirla: intervallo consigliato \\
$\lambda_w = 1$ & un'attivazione fuori fascia costa quanto lasciare quel consumo inspiegato \\
$\lambda_w > 1$ & la fascia diventa di fatto un vincolo rigido \\
\bottomrule
\end{tabularx}
\end{center}

\begin{warnbox}{Attenzione al valore predefinito}
Il valore predefinito $\lambda_w = 1$ colloca il termine \emph{esattamente} sul punto di pareggio:
accendere l'apparecchio fuori fascia costa quanto non spiegare affatto quel consumo. Il risultato
è che il modello, indifferente fra le due opzioni, tende a non attivare mai l'apparecchio fuori
orario --- la fascia si comporta come un vincolo rigido pur essendo formalmente flessibile.

Nelle verifiche sperimentali questo effetto è risultato marcato: con $\lambda_w = 1$ un
apparecchio non veniva attivato in nessun punto della giornata, dentro o fuori fascia. Per
osservare un comportamento realmente graduato è necessario scendere a $0.2$--$0.4$.
\end{warnbox}

\section{Termine 3 --- Numero di accensioni}

\begin{equation}
  \sum_{i} \sum_{t \in \Vset} \lambda_a \, p_i \, u_{i,t}
\end{equation}

\paragraph{Che cosa misura.} Il costo di ogni singolo avvio. Poiché $u_{i,t}$ vale 1 solo sui
fronti di salita, la somma $\sum_t u_{i,t}$ è esattamente il \textbf{numero di cicli} attribuiti
all'apparecchio nella giornata.

\paragraph{A che serve.} Senza questo termine il modello tende a produrre soluzioni frammentate:
accende e spegne ripetutamente gli apparecchi per inseguire ogni oscillazione del segnale,
generando decine di micro-cicli privi di senso fisico. Il termine introduce un costo fisso per
avvio, rendendo conveniente un unico ciclo lungo rispetto a molti brevi a parità di energia
attribuita.

\paragraph{Interpretazione del coefficiente.} Con $\lambda_a = 1$ un'accensione in più costa quanto
tenere l'apparecchio acceso per uno slot a fronte di segnale nullo. In altri termini: il modello
accetta un ciclo aggiuntivo solo se questo gli permette di spiegare almeno $p_i$ watt per uno slot
che altrimenti resterebbero inspiegati.

\paragraph{Comportamento al variare.} Valori crescenti producono soluzioni progressivamente più
``pulite'', con cicli meno numerosi e più lunghi. Valori troppo alti portano il modello a
rinunciare del tutto ad apparecchi che compaiono raramente ma realmente.

\section{Termine 4 --- Durata dei cicli}

\begin{equation}
  \sum_{i} \sum_{t} \lambda_d \, p_i \bigl( s^{-}_{i,t} + s^{+}_{i,t} \bigr)
\end{equation}

\paragraph{Che cosa misura.} Lo scostamento della durata di ogni ciclo dai limiti dichiarati. Le
due variabili sono asimmetriche nel significato ma simmetriche nel costo: $s^{-}$ misura di quanti
slot un ciclo è più \emph{corto} del minimo dichiarato, $s^{+}$ di quanti è più \emph{lungo} del
massimo.

\paragraph{Il punto non ovvio.} Come dimostrato nel paragrafo~\ref{sec:m4}, il costo complessivo di
un ciclo di lunghezza $L > b_i$ risulta
\begin{equation}
  \lambda_d \, p_i \, (L - b_i),
\end{equation}
cioè \textbf{esattamente proporzionale allo sforamento}. Questa proporzionalità non è imposta: è
una conseguenza del fatto che le finestre di vincolo violate sono tante quanti gli slot di
eccedenza. È il motivo per cui questa formulazione è preferibile a costruzioni più elaborate con
variabili di lunghezza del ciclo, che richiederebbero un numero di binarie proporzionale a
$T \times d_{\max}$.

\paragraph{Interpretazione del coefficiente.} Con $\lambda_d = 1$ ogni slot di sforamento costa
quanto uno slot di consumo inspiegato. Il modello estende un ciclo oltre il limite dichiarato solo
se ciò gli consente di spiegare almeno altrettanto segnale. Il valore predefinito $0.5$ rende lo
sforamento relativamente conveniente, coerentemente con la constatazione che le durate dichiarate
sono poco affidabili.

\section{Termine 5 --- Monte ore giornaliero}

\begin{equation}
  \sum_{i} \lambda_D \, p_i \bigl( d^{+}_i + d^{-}_i \bigr)
\end{equation}

\paragraph{Che cosa misura.} Lo scarto fra le ore complessive di accensione attribuite
all'apparecchio nella giornata e quelle attese in base a frequenza e durata dichiarate.

\paragraph{Differenza rispetto al termine precedente.} È un vincolo \textbf{globale} sulla
giornata: non guarda la forma dei cicli ma solo il loro totale. Un ciclo da quattro ore e quattro
cicli da un'ora sono per questo termine equivalenti. I due termini sono quindi complementari ---
uno controlla la forma, l'altro il volume --- e possono essere usati insieme.

\paragraph{Perché è disattivato per impostazione predefinita.} Sulle abitazioni del campione la sua
attivazione peggiorava leggermente l'errore di ricostruzione. L'attesa calcolata dal questionario
non corrispondeva all'uso reale osservato, e forzare il modello ad avvicinarvisi degradava la
qualità. Il valore predefinito è quindi $\lambda_D = 0$, con il termine disponibile per
riattivazione se il questionario dovesse risultare più affidabile.

\section{Termine 6 --- Quota settimanale}
\label{sec:t-quota}

\begin{equation}
  \sum_{i} \lambda_{q,i} \, p_i \, e_i
\end{equation}

\paragraph{Che cosa misura.} Il numero di accensioni attribuite nel giorno corrente che eccedono la
quota settimanale ancora disponibile.

\paragraph{L'eccezione alla regola del punto di pareggio.} Questo è \textbf{l'unico termine il cui
punto di pareggio non è 1}, e la ragione merita attenzione perché è la stessa che ha reso
necessaria una calibrazione per apparecchio.

Gli altri termini penalizzano quantità commisurate a \emph{uno slot}: uno slot fuori fascia, uno
slot di sforamento. Questo penalizza invece un'\emph{intera attivazione}, che dura tipicamente
$s_i$ slot. Il beneficio che il modello ottiene attribuendo un ciclo completo è quindi dell'ordine
di $s_i \cdot p_i$, e per contrastarlo il sovrapprezzo deve superare quel valore:
\begin{equation}
  \lambda_{q,i} = \code{over\_activation\_factor} \times s_i .
\end{equation}

\begin{notebox}{Verifica sperimentale del punto di pareggio}
Su un caso di prova con attivazioni da 4 slot, il comportamento del modello cambia
\textbf{esattamente} a $\lambda_q = 4$: sotto quel valore le accensioni eccedenti vengono comunque
attribuite, sopra vengono rifiutate. Il risultato coincide con la previsione analitica e conferma
la necessità di scalare il coefficiente per la durata tipica.

Un coefficiente scalare unico sarebbe mal calibrato: nel campione le durate tipiche vanno da 1
slot (microonde) a 8 slot (televisori), un fattore otto di differenza nel punto di pareggio.
\end{notebox}

\paragraph{Struttura del costo.} Il termine è nullo finché le accensioni restano entro la quota, e
cresce linearmente oltre. Non è quindi una penalità sull'uso dell'apparecchio in sé, ma
esclusivamente sul suo \emph{uso eccessivo} rispetto a quanto dichiarato.

\section{Termine 7 --- Stagionalità}

\begin{equation}
  \sum_{i} \sum_{\substack{t \in \Vset \\ \text{mese}(t) \notin \mathcal{M}_i}}
  \lambda_s \, p_i \, x_{i,t}
\end{equation}

\paragraph{Che cosa misura.} Il costo di attribuire consumo a un apparecchio in un mese in cui
l'utente ha dichiarato di non usarlo.

\paragraph{Perché il coefficiente è alto.} Il valore predefinito $\lambda_s = 8$ è nettamente
superiore al punto di pareggio, e questo è deliberato. La dichiarazione stagionale è
qualitativamente più affidabile delle altre: ricordare che il climatizzatore si usa solo d'estate
è molto più facile che ricordare a che ora si avvia di solito la lavatrice. Un valore alto rende
l'apparecchio praticamente assente fuori stagione.

\paragraph{Perché resta comunque flessibile.} Nonostante il valore elevato il termine è una
penalità e non un'esclusione: se il segnale contenesse un'evidenza schiacciante --- un plateau
prolungato che nessun altro apparecchio può spiegare --- il modello potrebbe ancora attribuirlo.
Un'esclusione rigida ($x_{i,t} = 0$ fuori stagione) sarebbe stata più veloce da risolvere, ma
avrebbe reso irrecuperabile un errore di compilazione del questionario.

\paragraph{Impatto osservato.} Su un'abitazione del campione il climatizzatore, dichiarato attivo
da giugno a settembre, riceveva $8.1\unit{kWh}$ nei blocchi di novembre, dicembre e gennaio.
Con il termine attivo scende a zero.

\section{Termine 8 --- Specificità della fascia oraria}

Questo termine non è additivo ma \textbf{moltiplicativo}: agisce come pavimento sul fattore
$f_i(t)$ del termine~2 e non introduce una sommatoria separata. La sua costruzione è nel
capitolo~\ref{ch:fattore}, dove viene trattata insieme al resto del fattore temporale.

\begin{notebox}{Perché una penalità e non un premio}
L'obiettivo di questo termine è favorire un apparecchio che ha dichiarato una fascia stretta
rispetto a uno che non ha dichiarato nulla. Si potrebbe ottenerlo in due modi: premiando il primo
o penalizzando il secondo. La scelta è caduta sulla penalità per una ragione precisa.

Un premio sarebbe un \textbf{coefficiente negativo} nella funzione obiettivo. Il solver, che
minimizza, avrebbe allora convenienza ad accendere apparecchi \emph{solo per incassare il premio},
anche in totale assenza di segnale da spiegare. Il fenomeno è controllabile mantenendo il premio
inferiore al costo di ricostruzione che ne deriverebbe, ma resta una fragilità strutturale del
modello.

La penalità sulla genericità produce \textbf{lo stesso ordinamento relativo} fra i due apparecchi
mantenendo tutti i coefficienti non negativi, e quindi nessun incentivo ad attivazioni spurie.
\end{notebox}
